% >>> The current consolidated user guide is docs/PUV_Pipeline_Guide.tex (built to
%     PUV_Pipeline_Guide.pdf). This file is retained as source material / deep reference;
%     its L1 quality-control description predates the 2026-07 per-channel QC rework.
%     Start newcomers on the guide.
% PUV Pipeline — Validation Results Section
% Drop into your paper's results section.

\subsection{PUV--MOP validation}

\subsubsection{Bulk wave parameter agreement}

Of 40 instruments across 20 deployments, 33 were processed successfully
(82.5\%), with 7 failures attributable to hardware issues (burial, bent
pipes, dead batteries). PUV-derived wave parameters showed strong
agreement with the MOP wave model across all four TBR23 instruments
(Table~\ref{tab:validation}). For instruments unaffected by biofouling,
$H_s$ correlations exceeded $R^2 = 0.83$ with RMSE $<
\SI{0.06}{\meter}$ and biases of \SIrange{+0.002}{+0.008}{\meter}. One
instrument (MOP580 5m) exhibited lower agreement ($R^2 = 0.66$, bias $=
\SI{+0.037}{\meter}$), attributed to intermittent kelp fouling that
reduced the valid segment rate to 49\% compared with
\SIrange{85}{91}{\percent} for the remaining instruments. A small but
consistent positive $H_s$ bias was observed across all instruments, with
PUV-measured wave heights systematically exceeding shoaled MOP values.

% TODO: Insert Table \ref{tab:validation} with R2, RMSE, bias, skill, N
% for Hs, Tm, Ef across all instruments

\subsubsection{Frequency-dependent spectral bias}

The frequency-resolved spectral ratio
$S_\eta^\text{PUV}(f) / S_\eta^\text{MOP,shoaled}(f)$, averaged across
all matched hours, revealed that the positive $H_s$ bias was concentrated
at the swell spectral peak (\SIrange{0.04}{0.08}{\hertz}), where the
pressure transfer function $K_p > 0.9$ and the correction is minimal
(Table~\ref{tab:spectral_ratio}). The mid-swell band
(\SIrange{0.08}{0.14}{\hertz}) showed near-unity ratios ($0.99$--$1.03$),
and the high-swell band (\SIrange{0.14}{0.25}{\hertz}), where the
pressure correction is most aggressive ($K_p$ as low as $0.24$), also
showed only modest excess ($1.03$--$1.05$). The low-swell amplification
increased with decreasing depth, from a band-averaged ratio of $1.17$ at
\SI{8.3}{\meter} to $1.28$ at \SI{5.5}{\meter}. The overall excess
energy relative to the shoaled MOP spectrum was \SI{4.5}{\percent} at
\SI{8.3}{\meter} depth and \SI{8.7}{\percent} at \SI{5.5}{\meter}.

The concentration of excess energy at the spectral peak---where the
pressure correction factor is close to unity---indicates that the bias
does not originate from the pressure-to-surface-elevation transfer
function.

% TODO: Insert Table \ref{tab:spectral_ratio} with band-averaged ratios
% and Kp ranges for each instrument

% Rewritten 2026-07-27. The earlier version of this section eliminated three
% mechanisms in favour of a fourth (MOP peak broadening) that was subsequently
% shown to be an artifact of comparing spectra on mismatched frequency grids.
% Re-tested on 72,948 hours / 61 records with
% validation/redo_hypothesis_elimination.m, the elimination inverted: nonlinear
% shoaling is the mechanism, bound long waves are its difference-frequency
% counterpart, directional narrowing is a real but separate bias, and peak
% broadening is eliminated. Full account:
% ../PUV_paper/docs/findings_hypothesis_elimination_2026-07-27.md
\subsubsection{Hypothesis testing}

Four candidate mechanisms for the systematic positive bias were evaluated
(Table~\ref{tab:hypotheses}).

\paragraph{Nonlinear shoaling.}
Triad interactions during shoaling transfer energy from the spectral peak to its
super-harmonic near $2f_p$. The PUV/model energy ratio, formed on the model's
own frequency bins and binned by $f/f_p$, shows that transfer directly: a trough
of $0.94$--$0.95$ at $f/f_p = 1.3$--$1.5$ and a maximum of $1.34$ at
$f/f_p = 2.1$ for hours with $H_s/h > 0.12$, against $1.10$ for hours below it.
The discrepancy is organised by the Ursell number
($\rho = +0.332$, $N = \num{72948}$ hours across 61 records) rather than by
depth, which drops to $\rho = -0.114$ once $H_s/h$ is held fixed. Binned by
$H_s/h$, the spectral-width ratio $\nu^\text{PUV}/\nu^\text{MOP}$ rises
monotonically from $1.000$ below $H_s/h = 0.04$ to $1.100$ above $0.20$: the
model and the observations agree to three decimal places when the waves are
linear, and diverge as they cease to be. Bicoherence at the swell
self-interaction orders the harmonic excess \emph{within} fixed $H_s/h$ bands,
from $\rho = -0.008$ below $H_s/h = 0.06$ to $\rho = +0.285$ above $0.20$ ---
flat where triad coupling cannot operate, strengthening where it can.

An earlier version of this analysis correlated the Ursell number against the
\emph{low-swell} band ratio (\SIrange{0.04}{0.08}{\hertz}) and found
$R^2 = 0.004$. That test cannot detect the mechanism: nonlinear shoaling deposits
energy near $2f_p \approx \SI{0.18}{\hertz}$, not at the peak.

\paragraph{Directional narrowing.}
The first-harmonic directional spread $\sigma_1$, from the PUV cross-spectral
$a_1$ and $b_1$, is \emph{wider} than the MOP spread at swell frequencies, by a
median factor of $1.156$ across 61 records ($p = \num{2.6e-6}$). The sign is
opposite to the narrowing expected from refraction-driven focusing, and the
wider PUV spread is consistent with scattered energy at shallow depths that the
forward-propagating model does not carry. The spread ratio does not covary with
the discrepancy it was proposed to explain ($p = 0.30$ against the $H_s$ ratio,
$p = 0.54$ against $\nu$), so it is a separate bias in the model's directional
distribution rather than the mechanism behind the energy difference. Because
$\sigma_1$ depends on $r_1 = \sqrt{a_1^2 + b_1^2}$, it is invariant to
instrument rotation and unaffected by the heading corrections applied in 2026-07.

\paragraph{Bound long waves.}
Second-order difference-frequency forcing predicts bound infragravity variance
from the swell band via the Hasselmann (1962) / Sch\"affer--Madsen (1995) kernel.
The ratio of predicted-bound to observed-total IG variance rises with $H_s/h$
($\rho = +0.913$ across 61 record medians; $+0.758$ per hour) and crosses unity
near $H_s/h = 0.12$, above which second-order theory over-predicts the bound
component. Against the spectral-width discrepancy the same ratio gives
$\rho = +0.288$, falling to $+0.149$ once $H_s/h$ is controlled. Bound-wave
forcing and super-harmonic generation both arise from triad interactions among
the same swell components, so this collinearity is expected; the two are aspects
of one nonlinear response rather than competing explanations, and the data here
do not separate them.

An earlier version rejected this hypothesis partly on the grounds that the excess
energy sat at the swell peak rather than at group frequencies. That apparent
concentration at the peak was itself the grid artifact described below.

% =====================================================================
% RETRACTED 2026-07-25 — READ BEFORE REUSING ANYTHING IN THIS PARAGRAPH
% OR THE "Cross-deployment confirmation" SUBSECTION THAT FOLLOWS.
%
% The peak-broadening result below is an ARTIFACT of the comparison
% method, not a property of the MOP model. It was computed by
% interpolating MOP's ~20 coarse frequency bins UP onto the PUV's
% 3601-point grid (validation/analyze_spectral_shape.m:110) and then
% measuring half-power bandwidth and Qp on the fine grid. MOP bandwidths of
% 0.005-0.10 Hz against a multi-taper resolution bandwidth of 2*NW/T =
% 2.2e-3 Hz is a ~2-45x mismatch. (Stated as "74x" before 2026-07-27, which
% compared against the BIN SPACING of 2.78e-4 Hz; an NW = 4 estimate is
% oversampled 8x relative to what it resolves. The empirical 68.3% artifact
% is unaffected -- it was measured, not derived from the ratio.) Linear interpolation between coarse bin centres
% manufactures a broad smooth peak and depresses int(S^2 df), the
% numerator of Qp. Both metrics move in the direction of the claim.
%
% Measured against a known-zero truth: a PUV spectrum compared against a
% degraded copy of ITSELF through that same path yields 68.3% apparent
% bandwidth narrowing and a Qp ratio of 1.225 (TBR23/MOP586_7m, 1382
% segments). That EXCEEDS the entire effect reported here. The m0 ratio
% of 0.9954 confirms it is a shape artifact, not an energy error.
%
% On a resolution-matched grid (both products reduced to the model's
% native bins first) the effect is absent and changes sign across sites:
% Qp ratio 1.093 Torrey, 1.041 Cardiff, 0.930 Coronado, 0.999 Torrey
% offshore. The matched metric recovers injected broadening at 95-96%
% efficiency down to r = 1.02, so this is an informative null.
%
% The factor-of-3 disagreement between the 16.7% here and the 36-78% in
% cross_deployment_shape_summary.mat is explained: the apparent effect
% scales with the PUV estimator's resolution (Welch/17-min vs
% multi-taper/1-hour), not with anything physical.
%
% Replacement code:  validation/compare_shape_matched.m
%                    validation/run_matched_shape_sweep.m
% Regression tests:  validation/test_resolution_artifact.m
%                    validation/test_shape_metric_sensitivity.m
%                    validation/test_shoal_bin_synthetic.m
% Full write-up:     PUV_paper/docs/findings_resolution_artifact_2026-07-24.md
%
% The hypothesis-elimination prose above was REWRITTEN on 2026-07-27 to reflect
% the redone analysis. This banner is retained because the paragraph it guards
% (the peak-broadening result) is still present as the record of what was
% originally claimed. See
% validation/redo_hypothesis_elimination.m and
% ../PUV_paper/docs/findings_hypothesis_elimination_2026-07-27.md:
%
%   H1 nonlinear shoaling    was "ruled out" -> IS THE MECHANISM
%       rho(Ursell,nu) = +0.332; nu = 1.000 in the most linear Hs/h bin;
%       excess localised at f/fp = 2.1 (1.338 at high Hs/h); bicoherence
%       orders it within Hs/h bands (+0.285 above Hs/h = 0.20)
%   H2 directional narrowing was "ruled out" -> REAL but SEPARATE bias
%       sigma_1 ratio 1.156, p = 2.6e-6, yet uncorrelated with the
%       discrepancy (p = 0.30, 0.54). Report it; it is not the cause.
%   H3 bound long waves      was "ruled out" -> REAL, collinear with H1
%       +0.913 per record vs Hs/h, but only +0.149 against nu after
%       controlling for Hs/h
%   H4 peak broadening       was THE WINNER  -> ELIMINATED (artifact)
%
% The three sections below that assert these were "ruled out" must be
% REWRITTEN, not annotated. The original analysis did not merely get the
% answer wrong -- it eliminated the correct one.
% =====================================================================

\paragraph{MOP spectral peak broadening.}
The Goda peakedness parameter $Q_p$ was systematically higher for PUV
spectra than for shoaled MOP spectra (median $Q_p = 1.88$ vs.\ $1.75$,
ratio $1.07$). The PUV half-power spectral bandwidth was
\SI{16.7}{\percent} narrower than MOP (\SI{0.039}{\hertz} vs.\
\SI{0.047}{\hertz}), and the PUV peak spectral density was
\SI{14.4}{\percent} higher. The peakedness difference
$\Delta Q_p = Q_p^\text{PUV} - Q_p^\text{MOP}$ showed the strongest
correlation with $H_s$ amplification of any tested mechanism
($R = 0.44$, $R^2 = 0.20$). Hours when the PUV spectrum was more peaked
relative to MOP corresponded to hours with larger positive $H_s$ bias.

These results indicate that the MOP spectral wave model introduces
numerical broadening of the spectral peak during propagation from the
initialization depth (\SI{10}{\meter}) to the instrument depth. The
in-situ PUV measurement captures the true, narrower spectral peak shape,
resulting in higher peak energy density and thus higher $H_s$ for a
given total variance. This spectral broadening is a recognized property
of discretized spectral wave models \citep[e.g.,][]{rogers2002,booij1999}
and has not been previously quantified using co-located PUV observations
in this depth range.

\subsubsection{Cross-deployment confirmation}

% RETRACTED 2026-07-25 along with the paragraph above -- see that banner.
% Two separate problems with this subsection:
%   (1) The pattern it reports is the resolution artifact, which is present
%       in every instrument by construction. Reproducibility across
%       deployments is exactly what an artifact produces; consistency of
%       sign was read as evidence of robustness when it was evidence of a
%       shared upstream computation.
%   (2) "11 instruments" is stale even on its own terms. The saved result
%       (outputs/validation/cross_deployment_shape_summary.mat) holds 33
%       instrument-records, and the current catalog is 65. The 11 figure
%       also appears in PIPELINE_NOTES.md and docs/puv_pipeline_slides.tex.

On a resolution-matched grid the peak-shape difference is absent across the
catalog: the median $Q_p$ ratio is $1.008$ over 62 instrument-records
($p = 0.44$), with an interquartile range of $0.980$--$1.039$ that spans unity.
The legacy interpolated-grid calculation returns $1.203$ on the same records,
within round-off of the $1.225$ obtained by comparing a PUV spectrum against a
degraded copy of itself.

Consistency of the original pattern across instruments was read as evidence that
it was physical. It is not: the artifact is a property of the comparison, so it
appears in every record by construction, and agreement across deployments at
matched depths is what a shared upstream error produces.

% TODO: Insert Table \ref{tab:hypotheses} summarizing hypothesis,
% test metric, result, and verdict for each of the four mechanisms

\subsubsection{Pipeline cross-validation: Ruby2D head-to-head}
\label{sec:ruby2d-validation}

To verify that the multi-taper L2 implementation does not introduce
biases relative to an established processing pipeline, we reprocessed
the Ruby2D MOP582\_6m record (Torrey Pines, October 2021 to February 2022,
$\sim$3,100 hours) and compared it segment-by-segment against the
archived L2 product produced by an independent prior pipeline
(\texttt{beachpeeps/PUV\_Processing}). The prior
pipeline uses a single-periodogram estimator (Hanning window, full
60-minute segments, $\Delta f = 0.000278$~\si{\hertz}) and applies a
fixed cap of $K_p^2 \le 10$ to the pressure transfer function. The legacy
record has 3,095 60-min segments; the new pipeline produces 10,903 17-min
segments over the same record. Segments were matched on midpoint
nearest neighbor with a tolerance of 30~minutes, yielding 2,322 matched
pairs (75\%). The remaining 25\% are hours that the new pipeline rejects
under the segment-validity criterion (\texttt{nanMaxFrac} = 0.10) but
that the prior pipeline retains by zero-padding small NaN gaps.

\begin{table}[h]
  \centering
  \caption{Ruby2D MOP582\_6m: bulk wave parameter agreement between
    the new multi-taper pipeline and the legacy archived L2. $T_p$ statistics
    are computed after restricting both estimates to the SS band
    (\SIrange{4}{25}{\second}); without this filter, three segments on
    January 15, 2022 in the prior pipeline report $T_p$ at the IG/DC
    bin ($T_p = 1200$, 1800, 3600~\si{\second}) because the prior
    pipeline's peak picker is not band-limited.}
  \label{tab:ruby2d}
  \begin{tabular}{@{}lccccc@{}}
    \toprule
    Metric & Legacy (med.) & Multi-taper (med.) & Bias (M$-$L) & RMS & $R^2$ \\
    \midrule
    $H_s$ (\si{\meter})        & 0.771  & 0.757  & $-0.009$ & 0.053 & 0.981 \\
    $T_p$ (\si{\second})       & 13.48  & 13.30  & $-0.02$  & 1.48  & 0.761 \\
    Direction SS (\si{\degree})    & $-7.3$ & $-7.2$ & $+0.2$   & 1.2   & 0.927 \\
    Spread SS (\si{\degree})       & 14.7   & 15.8   & $+1.0$   & 1.7   & 0.877 \\
    \bottomrule
  \end{tabular}
\end{table}

Bulk wave parameter agreement is excellent across all four metrics
(Table~\ref{tab:ruby2d}; Figure~\ref{fig:ruby2d_scatter}). $H_s$ agrees
to 5~cm RMS with a $-9$~mm median bias
($R^2 = 0.98$), and direction agrees to within
1.2$^\circ$~RMS ($R^2 = 0.93$). The directional spread is
1$^\circ$ wider in the new pipeline on average; this is the largest
systematic offset observed and remains within the segment-to-segment
scatter. The new $T_p$ shows the discrete quantization bands expected from
$\Delta f = 0.000976$~\si{\hertz} versus the prior pipeline's
$0.000278$~\si{\hertz}, but no systematic bias.

\begin{figure}[h]
  \centering
  \includegraphics[width=0.85\textwidth]{Ruby2D/ruby2d_582_6m_scatter.png}
  \caption{Bulk parameter agreement between the new multi-taper pipeline
    (vertical axis) and the legacy archived L2 (horizontal axis) for
    Ruby2D MOP582\_6m, 2,322 matched 60-minute segments,
    October 2021 to February 2022. The dashed line is 1:1.}
  \label{fig:ruby2d_scatter}
\end{figure}

Per-segment surface elevation spectra agree closely in shape
(Figure~\ref{fig:ruby2d_spectra}). At a representative median segment
($H_s = 0.77$~\si{\meter}, November 27, 2021) and at a storm peak
($H_s = 2.78$~\si{\meter}, December 14, 2021, bimodal sea state), the
integrated $H_s$ values from the two pipelines agree to within 4\%.
The visible difference between the spectra is the chi-square noise
characteristic of a single-periodogram estimator in the prior pipeline,
versus the variance reduction provided by averaging $K = 7$ DPSS tapers
in the new pipeline. The underlying peak shapes and locations agree.

\begin{figure}[h]
  \centering
  \includegraphics[width=0.95\textwidth]{Ruby2D/ruby2d_582_6m_spectra.png}
  \caption{Representative surface elevation spectra for Ruby2D
    MOP582\_6m. Top: median segment ($H_s = 0.77$~\si{\meter}). Bottom:
    storm peak ($H_s = 2.78$~\si{\meter}). Left column: log scale,
    full SS band. Right column: linear scale, swell band detail. The
    multi-taper estimates (blue) are smoother than the prior single-periodogram
    estimates (red); peak locations and total variance agree.}
  \label{fig:ruby2d_spectra}
\end{figure}

This cross-validation establishes two things. First, the multi-taper
implementation reproduces an independent prior pipeline to within
instrument noise on bulk wave parameters, confirming that the
estimator change is a refinement rather than a methodological break.
Prior published Ruby2D bulk wave statistics do not require revision.
Second, the consistency-of-windows fix in the new cross-spectral
calculations (auto- and cross-spectra both use the same DPSS tapers,
unlike the prior pipeline which mixed Hanning and Hamming) does not
produce the order-20\% directional errors that the bug analysis
suggested were possible at this site --- the practical impact at
Ruby2D MOP582\_6m is approximately $1^\circ$ in spread.

\subsubsection{Site-specific amplification at SIO Pier}

The SIO Pier instruments (MOP 511, ~7m depth) exhibited dramatically
larger PUV-to-MOP spectral amplification than all other sites. The
peak spectral density ratio ranged from 2.0 to 2.7 across four
deployment periods (SIO24A through SIO25E), with corresponding $H_s$
amplification of 29--36\% --- far exceeding the 1--11\% observed at
Torrey Pines, Solana Beach, and Los Pe\~nasquitos sites at comparable
depths.

Investigation confirmed that the MOP frequency grid is identical across
all stations (20 bins, 0.04--0.40~Hz), ruling out a resolution artifact.
The MOP model total energy at 10~m depth is comparable between SIO
($H_s = 0.62$~m) and Torrey Pines ($H_s = 0.61$~m), yet the in-situ
PUV at SIO measures approximately 2.5$\times$ the peak spectral density
predicted by the shoaled MOP spectrum.

The SIO Pier PUV is located near the head of Scripps Submarine Canyon,
where canyon-edge refraction is known to create wave energy focusing
\citep[e.g.,][]{munk1948,shepard1950}. The canyon bifurcates nearshore,
and wave focusing at the canyon head has been documented at Black's
Beach on the adjacent fork. The MOP model, which propagates spectra
along a relatively coarse alongshore grid, does not resolve this
fine-scale bathymetric focusing. The PUV measurements capture the true,
focused wave field, resulting in the observed amplification.

This site-specific effect is distinct from the universal spectral peak
broadening described above. While broadening produces a modest
($\sim$5--10\%) Hs bias that is consistent across all sites, the
canyon focusing at SIO produces an order-of-magnitude larger
amplification that is unique to that location.

\subsection{L3 wave forcing characterization}

\subsubsection{Energy band decomposition}

Frequency-band analysis revealed clear seasonal patterns across all
deployments. Winter deployments (TOR23W, SOL23, TOR24W, SOL24) were
swell-dominated (65--90\% of total energy flux in the 0.04--0.10~Hz
band), while summer deployments (TBR23, SIO25B--D) showed
approximately equal swell and sea contributions. The energy flux
consistency check confirmed exact agreement between the sum of band
fluxes and the total L2 energy flux ($F_\mathrm{total}/F_\mathrm{L2}
= 1.0000$) for all 33 instruments.

\subsubsection{Storm characterization with MOP gap-filling}

Storm detection identified 0 events during the calm summer TBR23
deployment and up to 11 events during winter TOR23W, including the
major December 28, 2023 storm sequence ($H_s$ up to 3.6~m at 7~m
depth, sustained over 74~hours). MOP gap-filling proved essential:
the SOL23 MOP651 7m instrument detected 0 storms from PUV data
(55\% valid segments) but MOP identified 3 events during data gaps,
including an 86-hour storm with $H_s = 3.09$~m. Across all
deployments, MOP supplemented 1--8 additional storm events per
instrument that were missed due to PUV data gaps during the most
energetic conditions.

\subsubsection{Transport proxy patterns}

The Shields parameter showed clear depth and seasonal dependence.
Median Shields numbers decreased from 0.24 at 5.6~m to 0.065 at
15.5~m depth (TOR23W MOP586 cross-shore array), with corresponding
mobilization fractions decreasing from 100\% to 78\%. Winter
deployments showed approximately 3$\times$ the cumulative bottom
energy flux of summer deployments at equivalent depths. The Rouse
number indicated bedload-dominated transport (median $P = 2.8$--6.0)
across all sites and seasons, with the lowest values (most
suspension-capable) at the shallowest winter sites.

\subsubsection{Tidal and subtidal currents}

Tidal currents were small (typically $<$0.05~m/s) relative to
wave-driven subtidal flows. The subtidal cross-shore residual was
consistently negative (offshore-directed) across all sites and
seasons, with magnitudes of 0.004--0.023~m/s, consistent with
wave-driven undertow. Tidal depth predictions from the Scripps Pier
NOAA gauge matched PUV depth anomalies at $R > 0.99$ for instruments
with continuous records, confirming the UTC clock convention and
validating the pressure-to-depth conversion.

\subsection{TBR23 post-storm recovery dynamics}

\subsubsection{Observed morphological change}

Jetski bathymetric surveys (22 surveys during the deployment period)
revealed 25--39~cm of net erosion at the PUV instrument locations
over the May--August 2023 deployment, with the 5~m sites eroding
more than the 7~m sites. Pre-deployment context from the MOPS survey
archive showed that the winter 2023 storm sequence had deposited
approximately 85~cm of sand at the 5~m depth contour relative to the
prior year, placing the bed well above its long-term equilibrium
profile. The TBR23 deployment thus captured the post-storm
readjustment phase rather than active storm erosion.

\subsubsection{Competing transport mechanisms}

Analysis of the PUV velocity moments revealed two competing
cross-shore transport mechanisms operating simultaneously during the
recovery period:

\paragraph{Onshore: velocity skewness.}
Velocity skewness $S_k$ was positive at all three reliable instruments
(median $S_k = +0.05$ to $+0.12$), indicating peaked crests and flat
troughs characteristic of shoaling waves. Skewness increased with
decreasing water depth, following the Ruessink et al.\ (2012)
parameterization based on the Ursell number. The skewness-driven
bedload flux $\langle u |u|^2 \rangle$ was directed onshore,
representing the wave-driven recovery mechanism.

\paragraph{Offshore: undertow.}
The subtidal mean cross-shore current was consistently offshore-directed
($\overline{u} = -0.006$ to $-0.012$~m/s) at all three instruments,
representing wave-driven return flow (undertow). Undertow magnitude
increased with wave height and was modulated by tidal phase, with
stronger offshore flow at high tide when more wave energy reached the
instrument depths.

\subsubsection{Depth-dependent transport balance}

The ratio of onshore skewness flux to offshore undertow flux showed a
clear depth dependence exploiting the tidal modulation of water depth
at each fixed instrument location ($\sim$2.5~m tidal range). At the
shallowest site (MOP586 5m, median depth 5.5~m), the skewness flux
exceeded the undertow flux by a factor of 3--4, indicating net onshore
transport and active recovery. At the 7~m sites (median depth 8.3~m),
the undertow flux slightly exceeded the skewness flux, indicating near-
equilibrium or weak offshore transport.

The crossover depth where onshore and offshore fluxes balanced
was approximately 6--7~m, defining a \emph{recovery boundary} below
which the post-storm deposit was actively being redistributed
shoreward. This is consistent with the survey observations: the
MOP586 5m site exhibited a 25~cm recovery pulse in mid-July 2023 that
was not observed at the deeper instruments.

One instrument (MOP580 5m) was excluded from the transport analysis
due to intermittent kelp fouling that produced anomalous negative
velocity skewness (92\% of segments) inconsistent with the other
instruments and with wave theory.

\subsection{Mean cross-shore flow validation}
\label{sec:mean_flow_results}

The validation framework defined in
Section~\ref{sec:mean_flow_validation} was applied across the
2023+ catalog (33 instruments at four sites: 19 at Torrey Pines,
4 at Los Pe\~nasquitos lagoon, 5 at Solana Beach, 5 at SIO Pier).
The same fits were repeated at the legacy 17-minute segmentation
to provide a segmentation-independence check; pre-2023 deployments
(Catalina Island, Ruby2D Torrey cross-shore array; total 5
additional instruments) were added at 1-hour only. The analysis
excludes the kelp-fouled TBR23 MOP580 5m record and the
SIO25C partial deployment ($N = 61$ segments, $R^2 = 0.001$).

\subsubsection{Cross-deployment headline}
The results decisively favor the real-Stokes-return-flow hypothesis
(Table~\ref{tab:mean_flow_headline}). The intercept distribution
is centered on zero across both segmentations (median $|\beta| =
\SI{0.6}{\milli\meter\per\second}$ at 17-min,
$\SI{2.8}{\milli\meter\per\second}$ at 1-hour), and 100\,\% of
17-min and 94\,\% of 1-hour intercepts fall inside a $\pm$2~cm/s
band that envelopes the manufacturer-quoted single-sample velocity
accuracy at the deployed range setting. The slope distribution
clusters near the Stokes prediction
($\alpha/\alpha_\text{theory} = +0.53$ at 17-min, $+0.71$ at
1-hour), with 66--81\,\% of instruments showing the predicted sign
of $\alpha < 0$ (offshore-directed undertow). Going from 17-min to
1-hour averaging \emph{strengthens} the slope-vs-theory match and
the correct-sign fraction; random instrument noise would do the
opposite.

% TODO: Insert Table \ref{tab:mean_flow_headline}
%   columns: metric | 17-min | 1-hour
%   rows: median |beta|; |beta|<2cm/s fraction; median alpha/alpha_th;
%         alpha<0 fraction; median R^2; median tidal modulation amplitude

\subsubsection{Site decomposition}
The Torrey Pines site (19 instruments across 5 deployments and
depths from 4.8 to 15.5~m) gives the cleanest theory match: 83\,\%
have the correct $\alpha$ sign and median $\alpha/\alpha_\text{theory}
= +0.77$ at 17-min. The Los Pe\~nasquitos lagoon site agrees
($+0.88$ ratio, 75\,\% correct sign at 17-min). At Solana Beach
and SIO Pier, the 17-min slopes are mixed in sign with low $R^2$,
but both sites \emph{flip} to predominantly negative $\alpha$ when
re-fit at 1-hour (Solana: 20\,\% $\to$ 100\,\% correct sign;
SIO Pier: 40\,\% $\to$ 80\,\%). This is interpreted as longer
averaging exposing a wave-driven signal that the 17-min sampling
variance was partially obscuring; the noise hypothesis predicts
the opposite shift.

\subsubsection{Tidal-phase modulation}
Across all instruments, the median peak-to-peak modulation
amplitude of $\overline{u}$ at $H_s \geq \SI{1}{\meter}$ is
\SI{2.6}{\centi\meter\per\second} (1-hour) and
\SI{2.8}{\centi\meter\per\second} (17-min), with peaks reaching
6--\SI{8}{\centi\meter\per\second} at multiple Torrey instruments.
A constant instrumental offset has zero modulation amplitude in
any $H_s$ stratum by definition; the observed amplitudes are 30--80
times the segment-mean random-noise floor and several times larger
than the largest fitted intercepts.

\subsubsection{Cross-instrument and Reynolds-stress consistency}
Time correlations of $\overline{u}(t)$ between independently-
calibrated Vectors at common timestamps were positive across all
non-fouled instrument pairs in TBR23, reaching $R = +0.57$ between
the two 7-m records. The kelp-fouled MOP580 5m record showed
\emph{anti}-correlation with all peers ($R = -0.33$ to $-0.48$),
providing an internal consistency check against the
fouling-flag exclusion. The Reynolds stress component
$\overline{u'w'}$, derived from the same record but via an
independent processing path, correlated positively with
$\overline{u}$ at all four TBR23 instruments
($R = +0.24$ to $+0.63$).

\subsubsection{Wave-direction discrimination}
The extended alongshore model
(Eq.~\ref{eq:wave_dir_check}) was decisive against the
shore-normal-rotation-error hypothesis. Across 30 instruments at
1-hour segmentation, the wave-direction-dependent coefficient
$\alpha_{v1}$ dominated the wave-direction-independent coefficient
$\alpha_{v0}$ in magnitude in 90\,\% of records, with the 95\,\%
confidence interval on $\alpha_{v1}$ excluding zero in 83\,\% of
records. The sign of $\alpha_{v1}$ was positive in 80\,\% of cases,
corresponding to NW swells driving southward longshore current and
S swells driving northward longshore current---the classical
radiation-stress signature for the WNW-facing San Diego County
coast. Reading $\alpha_{v0}/\alpha_u$ as an upper-bound implied
shore-normal angle error places that bound at a few degrees,
consistent with the precision of the CDIP MOP shore-normal
metadata.

\subsubsection{Verdict}
The cross-deployment dataset is incompatible with both the
constant-offset and the random-noise hypotheses, and matches the
Stokes return-flow prediction in sign, magnitude, and segmentation-
length response. Random Doppler noise reduces the segment-mean
noise floor to $\sim$0.4~mm/s at 17-min and $\sim$0.24~mm/s at
1-hour averaging, well below any of the slopes, modulation
amplitudes, or cross-instrument shared variance reported above. The
small (mm/s-scale) wave-direction-independent residuals visible in
the alongshore fit are consistent with site-specific Eulerian
residual currents and any small remaining shore-normal-rotation
error; they are not the dominant component of either velocity
component's response to wave forcing.
